\section{Class P --- Binary Arith --- Iterate}
%%% generated by the blueprint pipeline from TCSlib.Complexity.ClassP.BinaryArith.Iterate
\subsection*{Class P --- Binary Arith --- Iterate: Introduction}

Declarations of the Lean module \texttt{TCSlib.Complexity.ClassP.BinaryArith.Iterate}.

\subsection*{Class P --- Binary Arith --- Iterate: Results}

\begin{theorem}[Preservation under iteration]
\lean{Complexity.BinArith.iterate_invariant}
\label{Complexity.BinArith.iterate_invariant}
\leanok
\difficulty{4}
Let $\alpha$ be a type, let $g:\alpha\to\alpha$, and let $I$ be a property of elements
of $\alpha$. Suppose $I(s)$ implies $I(g(s))$ for every $s$. Then, for every natural
number $n$, if $I(s)$ holds, it also holds for the result of applying $g$ to $s$ exactly
$n$ times.
\end{theorem}

\begin{theorem}[Iteration reaches zero measure]
\lean{Complexity.BinArith.iterate_measure_eq_zero}
\label{Complexity.BinArith.iterate_measure_eq_zero}
\leanok
\difficulty{5}
Let $\alpha$ be a type, let $g:\alpha\to\alpha$, let $I$ be a property on $\alpha$, and
let $\mu:\alpha\to\mathbb{N}$. Suppose $I$ is preserved by $g$, that whenever $I(s)$ and
$0<\mu(s)$ hold the measure strictly decreases under $g$, and that whenever $I(s)$ and
$\mu(s)=0$ hold, $g(s)=s$. Then for every natural number $n$ and every $s$ satisfying
$I(s)$ and $\mu(s)\le n$, the measure after applying $g$ exactly $n$ times to $s$ is
zero.
\end{theorem}

\begin{definition}[Length-guarded state update]
\lean{Complexity.BinArith.guardNxt}
\label{Complexity.BinArith.guardNxt}
\leanok
\uses{Complexity.pairFstD, Complexity.pairSndD, Turing.pairEncode}
Let $g$ be a function from finite Boolean words to finite Boolean words, and let $h$ and
$r$ be Boolean words. The update discards $h$, takes the first component of the pair
encoded by $r$, and pairs it with either $g$ applied to the second component of $r$ (if
that output's length is at most the first component's length) or the second component of
$r$ itself.
\end{definition}

\begin{definition}[Guarded loop transformation]
\lean{Complexity.BinArith.guardLoop}
\label{Complexity.BinArith.guardLoop}
\leanok
\uses{Complexity.BinArith.guardNxt, Complexity.BlockLoop.step}
Let $g$ be a function from Boolean words to Boolean words. The guarded loop
transformation is the block-loop step whose state update applies the length-guarded
update and whose other argument is ignored.
\end{definition}

\begin{definition}[Bounds on encoded component lengths]
\lean{Complexity.BinArith.SlotsLe}
\label{Complexity.BinArith.SlotsLe}
\leanok
\uses{Complexity.pairFstD, Complexity.pairSndD}
For a natural number $L$ and a Boolean word $W$, write $a$ for the first component
decoded from $W$, $b$ for the first component decoded from the second component of $W$,
$c$ for the first component decoded from the second component of that second component,
and $d$ for the second component decoded from the second component of that second
component. The property holds exactly when each of $a,b,c,d$ has length at most $L$.
\end{definition}

\begin{theorem}[Self-bounded decoded slots]
\lean{Complexity.BinArith.slotsLe_self}
\label{Complexity.BinArith.slotsLe_self}
\leanok
\uses{Complexity.BinArith.SlotsLe, Complexity.length_pairFstD_le, Complexity.length_pairSndD_le, Complexity.pairSndD}
\difficulty{4}
For every finite bit string $W$, each of the four decoded slots—the first component of
$W$, the first component of its second component, the first component of the second
component of its second component, and the second component of the second component of
its second component—has length at most $|W|$.
\end{theorem}

\begin{theorem}[Two Possible Guard Loop Outcomes]
\lean{Complexity.BinArith.guardLoop_cases}
\label{Complexity.BinArith.guardLoop_cases}
\leanok
\uses{Complexity.BinArith.guardLoop, Complexity.BinArith.guardNxt, Complexity.BlockLoop.step, Complexity.blockDone, Complexity.pairFstD, Complexity.pairSndD, Turing.pairEncode}
\difficulty{4}
Let $g$ be a function on bit strings and let $W$ be a bit string. Write $\langle
a,b\rangle$ for the pairing of bit strings, and let $\pi_1$ and $\pi_2$ be the total
projections that return the empty string when parsing fails. The result of the guard
loop on $W$ is either the fixed completed-block encoding, or the encoding whose first
component is $\pi_1(W)$ and whose second component pairs the tail of $\pi_1(\pi_2(W))$
with the result of applying the guarded next-state function to $\pi_1(W)$ and
$\pi_2(\pi_2(W))$.
\end{theorem}

\begin{theorem}[Preservation of slot length bounds]
\lean{Complexity.BinArith.slotsLe_guardLoop}
\label{Complexity.BinArith.slotsLe_guardLoop}
\leanok
\uses{Complexity.BinArith.SlotsLe, Complexity.BinArith.guardLoop, Complexity.BinArith.guardLoop_cases, Complexity.BinArith.guardNxt, Complexity.blockDone, Complexity.pairFstD_nil, Complexity.pairFstD_pairEncode, Complexity.pairSndD, Complexity.pairSndD_pairEncode}
\difficulty{5}
Let $g$ be any function from finite bit strings to finite bit strings. For a finite bit
string $W$ and natural number $L$, suppose the four decoded slots of $W$ have lengths at
most $L$: the first component of $W$, the first component of its second component, the
first component of the second component of its second component, and the second
component of the second component of its second component. After one guarded loop
step—using $g$ only when its output length is at most the current first-slot length—the
four slots still have lengths at most $L$.
\end{theorem}

\begin{theorem}[Length bound for one guarded update]
\lean{Complexity.BinArith.length_guardLoop_le}
\label{Complexity.BinArith.length_guardLoop_le}
\leanok
\uses{Complexity.BinArith.SlotsLe, Complexity.BinArith.guardLoop, Complexity.BinArith.guardLoop_cases, Complexity.BinArith.guardNxt, Complexity.BinArith.slotsLe_guardLoop, Complexity.blockDone, Complexity.pairFstD_pairEncode, Complexity.pairSndD_pairEncode, Turing.length_pairEncode}
\difficulty{4}
Let $g$ be a function from finite bit strings to finite bit strings, and let $L$ be a
natural number. Given a bit string $W$, repeatedly decode its first and second
components, treating an unsuccessful parse as having empty components. Suppose the first
component of $W$, the first component of its second component, the first component of
its iterated second component, and the iterated second component three times each have
length at most $L$. Apply the guarded update that returns the encoding of two empty
words if the first component of $W$ or the first component of its second component is
empty; otherwise, it encodes the unchanged first component together with a nested
encoding whose first component is the first component of the second component with its
initial bit removed, and whose second component is the encoding of that first component
paired with the result of applying $g$ to the second component of the iterated second
component, provided the result has length at most the first component's length, and
otherwise to that second component itself. The resulting bit string has length at most
$7L+6$.
\end{theorem}

\begin{theorem}[Length bound for iterated guarded updates]
\lean{Complexity.BinArith.length_iterate_guardLoop_le}
\label{Complexity.BinArith.length_iterate_guardLoop_le}
\leanok
\uses{Complexity.BinArith.SlotsLe, Complexity.BinArith.guardLoop, Complexity.BinArith.iterate_invariant, Complexity.BinArith.length_guardLoop_le, Complexity.BinArith.slotsLe_guardLoop, Complexity.BinArith.slotsLe_self}
\difficulty{4}
For any function $g$ from finite bit strings to finite bit strings, any finite bit
string $w$, and any natural number $i$, the result of applying the guarded loop
determined by $g$ exactly $i$ times to $w$ has length at most $7(|w|+1)$.
\end{theorem}

\begin{theorem}[Polynomial-Time Computability of a Guarded Loop]
\lean{Complexity.BinArith.polyTimeComputable_guardLoop}
\label{Complexity.BinArith.polyTimeComputable_guardLoop}
\leanok
\uses{Complexity.BinArith.guardLoop, Complexity.BinArith.guardNxt, Complexity.BlockLoop.polyTimeComputable_step, Complexity.PolyTimeComputable, Complexity.PolyTimeComputable.comp, Complexity.PolyTimeComputable.pairEncode, Complexity.pairFstD, Complexity.pairSndD, Complexity.polyTimeComputable_ite, Complexity.polyTimeComputable_lenLe, Complexity.polyTimeComputable_pairFstD, Complexity.polyTimeComputable_pairSndD}
\difficulty{4}
Let $g$ be a polynomial-time computable function on bit strings. Define a loop state
update as follows: if either the first decoded component of the state is empty or the
first decoded component of its second decoded component is empty, replace the state by
the fixed encoding of two empty strings. Otherwise, keep the first component unchanged
and replace the second component by an encoding whose first component is obtained by
applying $g$ to the old second component's second decoded component when that image has
length at most the old first component's length, and is unchanged otherwise; the second
component of this inner encoding is the image under $g$ in the first case and the
unchanged old value in the second case. Then the function that applies this update once
is computable in polynomial time by a finite Turing machine.
\end{theorem}

\begin{theorem}[Polynomial-Time Bounded Iteration]
\lean{Complexity.BinArith.polyTimeComputable_iterate}
\label{Complexity.BinArith.polyTimeComputable_iterate}
\leanok
\uses{Complexity.BinArith.guardLoop, Complexity.BinArith.guardNxt, Complexity.BinArith.length_iterate_guardLoop_le, Complexity.BinArith.polyTimeComputable_guardLoop, Complexity.BlockLoop.loop_output, Complexity.BlockLoop.polyTimeComputable_countdownInit, Complexity.BlockLoop.polyTimeComputable_emit, Complexity.BlockLoop.polyTimeComputable_of_loop, Complexity.BlockLoop.step, Complexity.PolyTimeComputable, Complexity.PolyTimeComputable.comp, Complexity.PolyTimeComputable.pairEncode, Complexity.pairFstD_pairEncode, Complexity.pairSndD, Complexity.pairSndD_pairEncode, Complexity.polyTimeComputable_id, Complexity.polyTimeComputable_pairSndD, Complexity.polyTimeComputable_polyUnary, Turing.length_pairEncode, Turing.pairEncode}
\difficulty{6}
Let $g$ and $u$ be polynomial-time computable functions from finite Boolean words to
finite Boolean words. Let $a,k,B,b$ be natural numbers, and suppose that for every word
$z$ and every natural number $i$ with $i\le a(|z|+1)^k$, the word obtained by applying
$g$ successively $i$ times to $u(z)$ has length at most $B(|z|+1)^b$. Then the function
sending $z$ to the result of applying $g$ successively $a(|z|+1)^k$ times to $u(z)$ is
polynomial-time computable.
\end{theorem}
